\section{Risk Quantification}

\subsection{Dynamic Underwriting}

Every loan quote must bound the lender's loss before capital is deployed. The risk engine evaluates each collateral asset dynamically, deciding whether to accept it, reject it, or accept it with adjusted parameters, and continuously re-evaluates existing positions. Inputs include liquidity conditions (depth and slippage at trade size), non-quantitative dimensions (smart contract, issuer, and infrastructure risk), and venue conditions for Prime Trade. Maximum borrow, fixed rate, and liquidation diagnostics are \emph{outputs of the model}, not static protocol parameters.

\subsection{Methodology}

The risk engine follows a three-stage pipeline: data collection, statistical modeling, and simulation.

\textbf{Data collection.} The engine ingests live on-chain state (asset prices, lending-venue utilization rates, liquidity-pool depth, and historical return series) directly from the chains where collateral and principal reside.

\textbf{Statistical modeling.} Collected data feeds statistical models that estimate the current volatility regime, tail-risk characteristics, and correlation structure of the collateral asset. The models capture asymmetric behavior (downside moves weigh more than equivalent upside moves) and discontinuous price events that continuous models would understate.

\textbf{Simulation.} Monte Carlo simulations project collateral prices and variable-rate funding costs over the loan horizon under the estimated regime. Simulated paths are evaluated against the loan's liquidation barriers and actual liquidation mechanics (partial versus full close, bonuses, route slippage), so pricing reflects path-dependent recovery. The output --- a joint distribution of liquidation probability and loss given default --- feeds directly into the pricing module.

\subsection{Pricing}

The published credit multiplier and rate are chosen jointly to balance the lender's expected return, position economics, and incremental tail risk when the new loan sits beside existing positions.

Higher volatility or longer duration lowers the credit multiplier; adding correlated exposure to a concentrated book raises the rate. Positions that reduce aggregate tail risk receive more favorable terms, naturally incentivizing diversification. The rate never falls below the risk-free return plus the simulated expected loss per dollar lent. With principal $L$, annualized rate $r$, and horizon $T$ in years, debt at maturity is $D = L(1 + rT)$, linking the quoted rate directly to the liquidation barrier and simulated loss distribution.

For externally-sourced loans, the rate buffer is sized separately from the borrower's fixed rate, covering the tail of simulated variable-rate paths on the external venue. The borrower's obligation remains the simple linear formula above regardless of the sourcing path.
